\cotefile{77}
\scanpage{34}
L'ens des sections de $2F$ est \deleted{\unsure{vide ou}} \deleted{\unsure{(qui exprime|(qui est}} 
\deleted{torseur sous $\Gamma(S,L\unsure{)}$} non vide \unsure{sss|ssi} la classe $2\,\mathrm{cl}(F)=\mathrm{cl}(2F)$
dans $H^1(S,L)$ est $=0$. Alors cet ens est
un torseur sous $H^0(S,L)$ (qui exprime
la première indétermination). La section
$\underline b$ étant choisie, l'ens des $\sigma$ au dessus
de $\underline b$ est non vide sss la classe
\[
 c(\underline b)=\mathrm{cl}(\Sigma(\underline b))\in H^1(S,L^{\otimes2})
\]
est $=0$; et dans cet cas est un torseur
sous $H^{\unsure{1|0}}(S,L^{\otimes2})$.

On peut considérer que $\Sigma^2(F)$, torseur
affine sous $\check B\otimes L$, est donné par une
structure d'extension
\[
 \unsure{\cdot|0}\unsure{-|\to}\check B\otimes L\longrightarrow V^{(3)}\longrightarrow\mathcal O_S\longrightarrow0 \tag*{(45)}
\]
(où on considère $\check B$ comme $V^{(2)}$).
Or $\mathrm{Sym}^2(\check B)=\mathrm{Sym}^2(V^{(1)})$ est une
extension de $\mathrm{Sym}^2(\mathcal O_S)=\mathcal O_S$ par $\mathcal O_S\otimes V^{\unsure{(2)|(1)}}=V^{(1)}$.
En fait, la proposition précédente se
précise en
\[
\begin{aligned}
 \Sigma^2(F)={}&\text{fibré affine associé à}\\
 &\text{l'extension }V^{(3)}=\mathrm{Sym}^2(\check B)\\
 &\text{de }\mathcal O_S\text{ par }V^{\unsure{(2)|(2)'|(1)}}=\check B.
\end{aligned}
\tag*{(46)}
\]
% archive sheet number: 33
